\begin{problem}[第37届全国中学生物理竞赛决赛][脉冲星信号与星际介质]{129}
    在恒星之间的广阔宇宙空间中存在星际介质。在宇宙射线的作用下，星际介质中的分子失去部分电子成为正离子，电子则游离在外，成为等离子态。脉冲星是高速旋转、具有超强磁场的中子星，地球上所观测到其发出的电磁辐射是脉冲信号。脉冲星信号可用于星际导航和高精度计时，为此需要获得脉冲信号到达地球的精确时间，研究其电磁辐射的相位变化与色散。
    一脉冲星到地球的距离为$d$，星际介质中的电子平均数密度为$n _ { \mathrm { e } }$（数量级为$10^{4} \mathrm{m}^{-3}$）；假设介质中存在匀强静磁场$B$（数量级为$10^{-10} \mathrm{T}$），其方向平行于电磁波的传播方向。已知电子质量为$m _ { \mathrm { e } }$，电荷为$- e \ ( e > 0 )$，真空介电常量为$\varepsilon _ { 0 }$，真空中光速为$c$。
    \begin{subproblem}
        取该脉冲星到地球的电磁辐射方向为$z$轴正向，对于频率为$f$的电磁波，其电场为$E _ { x } = E _ { 0 } \cos ( k z - 2 \pi f t )$，$E _ { y } = E _ { 0 } \cos ( k z - 2 \pi f t \pm \frac { \pi } { 2 } )$，式中$E _ { 0 }$和$k$分别为电磁波的振幅和波数。为简单起见，设$E _ { 0 }$为常量、且电磁波本身的磁场对电子的作用可忽略，求电子运动的回旋半径$R _ { \mathrm { e } }$。
    \end{subproblem}
    \begin{subproblem}
        脉冲星信号在星际介质中传播时会发生色散，其传播速度大小（群速度的大小$v _ { \mathrm { g } }$）依赖于电磁波频率相对于波数的变化率：$v _ { \mathrm { g } } = 2 \pi \frac { \mathrm { d } f } { \mathrm { d } k }$。求能通过介质到达地球的脉冲星信号电磁波的最低频率$f _ { \mathrm { c } }$。
    \end{subproblem}
    \begin{subproblem}
        假设在脉冲星信号的传播路径上，星际介质中的电子平均数密度$n _ { \mathrm { e } }$保持不变，问脉冲星信号中频率为$f$的电磁波到达地球的时间比其在真空中传播的时间延迟了多少？
    \end{subproblem}
    \begin{subproblem}
        观测发现，从脉冲星同时出发的频率为$f$的电磁波到达地球时出现了相位差，求该相位差（可取近似：$\frac { e } { 2 \pi } \sqrt { \frac { n _ { \mathrm { e } } } { \varepsilon _ { 0 } m _ { \mathrm { e } } } } << f , \frac { e B } { 2 \pi m _ { \mathrm { e } } } << f$）。
    \end{subproblem}
    \begin{subproblem}
        如果在传播途中长度为$a$的区间内电子数密度出现涨落$\Delta n _ { \mathrm { e } }$，求脉冲星信号中频率为$f$的电磁波通过该区间后由电子数密度涨落引起的相位移动。
    \end{subproblem}
\end{problem}